\section{Semantic Metric and Embedding Theory}\label{sec:semantics}

\S\ref{sec:topology} built the chain complex and its homology~$H_k$ on each stage layer~$\Delta^{(m)}$.
We now add two layers of structure.
First, \emph{cohomology}~(\S\ref{sec:cohomology}): the dual theory that detects obstructions to consistent assignment. This requires no additional assumptions.
Second, \emph{semantic geometry}~(\S\ref{sec:semantic-hash}--\S\ref{sec:ai-frontier}): we add an LLM embedding on vertices and extend it to every stage via the \emph{stage embedding tower}. The tower yields a per-stage embedding obstruction theorem, a grade decomposition of distance, and formal characterizations of hallucination, scaling saturation, understanding, and chain-of-thought reasoning.


% ============================================================
%   §4.1  COHOMOLOGY
% ============================================================
\subsection{Cohomology}\label{sec:cohomology}

Homology detects ``holes'': cycles that do not bound. Cohomology is the dual theory: it detects obstructions to extending local assignments globally. The definitions in this subsection require no constraint and no embedding; they hold for any astrolabe file.

\begin{definition}[Cochain complex]
\label{def:cochain}
The $k$-th cochain group is $C^k(\Delta) = \operatorname{Hom}(C_k, \mathbb{Z})$, the group of $\mathbb{Z}$-valued functions on $k$-simplices. The coboundary operator $\delta^k: C^k \to C^{k+1}$ is the dual of~$\partial_{k+1}$:
\[
  (\delta^k f)(\sigma) = f(\partial_{k+1}\sigma), \quad f \in C^k, \; \sigma \in C_{k+1}.
\]
\end{definition}

\begin{definition}[Cohomology]
\label{def:cohomology}
The $k$-th cohomology group is
\[
  H^k(\Delta) = \ker \delta^k / \operatorname{im} \delta^{k-1}.
\]
\end{definition}

The cochain complex runs in the opposite direction to the chain complex:
\[
  0 \to C^0(\Delta) \xrightarrow{\delta^0} C^1(\Delta) \xrightarrow{\delta^1} C^2(\Delta) \to \cdots
\]
While homology detects ``where the holes are,'' cohomology detects obstructions to consistent assignment. A $k$-cochain assigns a value (a score, a label, a priority) to each $k$-simplex; cohomology measures whether such assignments can be made consistently.

\begin{itemize}
\item $H^0$: the number of independent ways to assign a constant label to all vertices within each connected component. If $\beta_0 = 3$, then $\operatorname{rank} H^0 = 3$.
\item $H^1$: non-conservative flows. Assign a ``cost'' to each edge (e.g., estimated review time). If these costs cannot be expressed as differences of vertex potentials ($\operatorname{cost}(e) \neq \phi(t) - \phi(s)$), the obstruction lives in~$H^1$.
\item $H^2$: irreducible higher-order effects. Obstructions to decomposing a ternary relationship into pairwise binary ones.
\end{itemize}

\begin{example}[$H^1$: code review cost]
A software project has four modules $A, B, C, D$ and edges $A \to B \to C \to D \to A$ with estimated review times $3, 2, 4, 1$ (hours). If this were conservative, vertex potentials~$\phi$ would satisfy $\phi(B)-\phi(A) = 3$, $\phi(C)-\phi(B) = 2$, $\phi(D)-\phi(C) = 4$, $\phi(A)-\phi(D) = 1$. But $3+2+4+1 = 10 \neq 0$: no consistent potential exists. This obstruction is a class in~$H^1$; its value measures the ``effort debt'' per cycle traversal.
\end{example}

\begin{example}[$H^2$: irreducible ternary interaction]
Three mathematical concepts (continuity, compactness, convergence) each have pairwise relationships. But the Heine--Cantor theorem (``continuous functions on compact sets are uniformly continuous'') is a ternary interaction that cannot be decomposed into three independent binary statements. If this ternary record is absent, the obstruction lives in~$H^2$.
\end{example}

\begin{example}[Cochain as annotation]
A $0$-cochain $f \in C^0$ assigns a numerical score to each vertex: $f(\mathrm{TP53}) = 8$, $f(\mathrm{MDM2}) = 5$. The coboundary $\delta^0 f$ is a $1$-cochain: $(\delta^0 f)(e) = f(\mathrm{MDM2}) - f(\mathrm{TP53}) = -3$. If a $1$-cochain $g$ equals~$\delta^0 f$, then edge scores are fully determined by vertex scores. If not, $g$ represents a nontrivial class in~$H^1$: edge-level information cannot be reduced to vertex-level information.
\end{example}

The Kronecker pairing and cup product are introduced in \S\ref{sec:obstruction} and \S\ref{sec:ai-frontier}, where they are used.


% ============================================================
%   §4.2  CLASSICAL CONSTRAINT AND SEMANTIC HASH
% ============================================================
\subsection{Classical Constraint and Semantic Hash}\label{sec:semantic-hash}

\begin{definition}[Classical constraint]\label{def:classical}
An astrolabe file~$\mathbb{A}$ satisfies the \emph{classical constraint} if every hash in $\texttt{ref}(\sigma)$ belongs to~$\mathcal{H}_0$ for all $\sigma \in \mathbb{A}$. Equivalently, every reference points to an atom. Under this constraint, the stage decomposition collapses: $\mathcal{S}_1 = \mathbb{A}$ and $\Delta^{(0)}$ is the only nonempty stage layer.
\end{definition}

\begin{definition}[Semantic hash]\label{def:semantic-hash}
A \emph{semantic hash} is a function $h_{\mathrm{emb}}: \mathcal{H}_0 \to \mathbb{R}^d$ that maps each atom to a point in $d$-dimensional Euclidean space, such that semantically similar records are mapped to nearby points. In practice, $h_{\mathrm{emb}}$ is computed by a language model embedding (e.g., OpenAI text-embedding-3, $d=3072$; Sentence-BERT, $d=768$).
\end{definition}

Only atoms require LLM calls. All higher-dimensional embeddings are computed automatically from atom embeddings via the stage embedding tower (Definition~\ref{def:stage-tower}).

\begin{definition}[Semantic metric]
\label{def:semantic-metric}
The semantic hash induces a metric on~$\mathcal{H}_0$:
\[
  d_{\mathrm{emb}}(v, v') = \|h_{\mathrm{emb}}(v) - h_{\mathrm{emb}}(v')\|_2.
\]
\end{definition}

The identity hash~$h_{\mathrm{id}}$ captures textual identity; the semantic hash~$h_{\mathrm{emb}}$ captures meaning similarity. Their interplay reveals two fundamental failure modes.

\begin{example}[Two faces of the $h_{\mathrm{id}}/h_{\mathrm{emb}}$ gap]\label{ex:poincare-pitts}
\emph{Same meaning, different text.} The Poincar\'{e} conjecture can be stated as $r_1 = ((\texttt{stmt}, \text{``simply connected closed 3-mfd} \cong S^3\text{''}))$ and $r_2 = ((\texttt{stmt}, \text{``compact 3-mfd with } \pi_1 = 0 \text{ is a sphere''}))$. These are the same theorem, but $h_{\mathrm{id}}$ assigns unrelated hashes.

This has a precise counterpart in type theory: $\beta$-, $\delta$-, $\eta$-, and $\iota$-reduction transform a term into a definitionally equal term with different syntax. In Lean~4, the kernel considers $\beta\delta\eta\iota$-equivalent terms identical, but $h_{\mathrm{id}}$ distinguishes them. Homotopy type theory goes further: the univalence axiom identifies equivalent types with equal types.

\emph{Same text, different meaning.} Changing a single digit in a regularity statement (``smooth for $n \leq 6$'' vs.\ ``smooth for $n \leq 7$'') turns a theorem into a false statement. Both $h_{\mathrm{id}}$ and $d_{\mathrm{emb}}$ see them as nearly identical; the proof dependency structure distinguishes them.
\end{example}


% ============================================================
%   §4.3  STAGE EMBEDDING TOWER
% ============================================================
\subsection{Stage Embedding Tower}\label{sec:stage-tower}

In \S\ref{sec:semantic-hash}, we defined $h_{\mathrm{emb}}: \mathcal{H}_0 \to \mathbb{R}^d$. This embeds atoms. But the vertex set of $\Delta^{(1)}$ is $\mathcal{H}_1$, which includes the edges and faces of $\Delta^{(0)}$; the vertex set of $\Delta^{(2)}$ is $\mathcal{H}_2$, which includes the simplices of $\Delta^{(1)}$; and so on. We need an embedding at every stage.

The key observation is that a $p$-simplex already has a natural embedding: its exterior product $\Phi_0(\sigma) \in \bigwedge^p \mathbb{R}^d$. When $\sigma$ becomes a vertex in the next stage, we use this element as its position. The ambient space for stage~$m+1$ is therefore the full exterior algebra of stage~$m$.

\begin{definition}[Stage embedding tower]\label{def:stage-tower}
The \emph{stage embedding tower} is the sequence of vector spaces, vertex embeddings, and exterior embeddings
\[
  (W_0, \varphi_0, \Phi_0), \quad (W_1, \varphi_1, \Phi_1), \quad (W_2, \varphi_2, \Phi_2), \quad \ldots
\]
defined recursively.

\textbf{Level~0.} $W_0 = \mathbb{R}^d$. The vertex embedding is $\varphi_0 = h_{\mathrm{emb}}: \mathcal{H}_0 \to W_0$. For $\sigma \in \Delta_k^{(0)}$ with $\mathrm{ref} = [v_0, \ldots, v_k]$, the exterior embedding is
\[
  \Phi_0(\sigma) = (\varphi_0(v_1) - \varphi_0(v_0)) \wedge \cdots \wedge (\varphi_0(v_k) - \varphi_0(v_0)) \in \textstyle\bigwedge^k W_0.
\]

\textbf{Level~$m+1$.} The ambient space is the full exterior algebra of the previous level:
\[
  W_{m+1} = \textstyle\bigwedge W_m = \bigoplus_{p=0}^{\dim W_m} \textstyle\bigwedge^p W_m.
\]
The vertex embedding $\varphi_{m+1}: \mathcal{H}_{m+1} \to W_{m+1}$ is defined as follows. Each $v \in \mathcal{H}_{m+1}$ corresponds to an entry~$\sigma_v$ that entered the stage decomposition at some stage layer $\Delta^{(s)}$, $s \leq m$, with degree $p = \deg(\sigma_v)$.
\begin{itemize}
\item If $\sigma_v$ \textbf{entered at stage~$m$} (i.e.\ $\sigma_v \in \Delta^{(m)}$) as a $p$-simplex:
\[
  \varphi_{m+1}(v) = \Phi_m(\sigma_v) \in \textstyle\bigwedge^p W_m \subset W_{m+1}.
\]
\item If $\sigma_v$ \textbf{entered at an earlier stage} $s < m$:
\[
  \varphi_{m+1}(v) = \iota(\varphi_m(v)),
\]
where $\iota: W_m \hookrightarrow \bigwedge^1 W_m \subset W_{m+1}$ is the canonical inclusion.
\end{itemize}
For $\sigma \in \Delta_k^{(m+1)}$ with $\mathrm{ref} = [w_0, \ldots, w_k]$, $w_i \in \mathcal{H}_{m+1}$, the exterior embedding is
\[
  \Phi_{m+1}(\sigma) = (\varphi_{m+1}(w_1) - \varphi_{m+1}(w_0)) \wedge \cdots \wedge (\varphi_{m+1}(w_k) - \varphi_{m+1}(w_0)) \in \textstyle\bigwedge^k W_{m+1}.
\]
\end{definition}

The tower is not a design choice but a forced construction. Stage decomposition says ``the next stage treats the previous stage's simplices as vertices.'' A $p$-simplex at stage $m$ has a natural position: its exterior product $\Phi_m(\sigma) \in \bigwedge^p W_m$. To accommodate all degrees simultaneously, the ambient space must contain $\bigoplus_p \bigwedge^p W_m = \bigwedge W_m$. This is why $W_{m+1} = \bigwedge W_m$: it is the smallest space that can serve as the vertex space for the next stage.

The price is super-exponential dimension growth: $d, 2^d, 2^{2^d}, \ldots$ In practice this is not a limitation but an asset: higher stages have exponentially more room to embed, and the binding constraint always comes from stage~0 (Theorem~\ref{thm:binding}).

The key identity connecting adjacent stages is: \emph{stage $m$ volume = stage $m+1$ position norm.} The exterior embedding $\Phi_m(\sigma)$ measures the $k$-volume of $\sigma$ at stage $m$; the same element, reinterpreted as $\varphi_{m+1}(v)$, is the position of $\sigma$ as a vertex at stage $m+1$. A simplex with large volume occupies a prominent position in the next stage; one with zero volume is invisible there.

\begin{definition}[Stage $k$-volume]
The $k$-volume of $\sigma \in \Delta_k^{(m)}$ is
\[
  \operatorname{vol}_k^{(m)}(\sigma) = \|\Phi_m(\sigma)\|,
\]
where the norm on $\bigwedge^k W_m$ is induced by the inner product on~$W_m$ (making different grades orthogonal).
\end{definition}

We now verify the basic properties of the tower: dimension growth, well-definedness, and compatibility with the stage structure.

\begin{proposition}[Dimension tower]\label{prop:dim-tower}
$\dim W_0 = d$ and $\dim W_{m+1} = 2^{\dim W_m}$. Explicitly:
\[
  \dim W_0 = d, \quad \dim W_1 = 2^d, \quad \dim W_2 = 2^{2^d}, \quad \dim W_m = \underbrace{2^{2^{\cdot^{\cdot^{\cdot^d}}}}}_{m \text{ exponentials}}.
\]
\end{proposition}

\begin{proof}
$\dim \bigwedge V = 2^{\dim V}$ for any finite-dimensional vector space~$V$.
\end{proof}

\begin{proposition}[Well-definedness]
$\varphi_{m+1}$ is well-defined: every vertex in $\mathcal{H}_{m+1}$ receives exactly one embedding in~$W_{m+1}$.
\end{proposition}

\begin{proof}
Each entry~$\sigma_v$ enters the stage decomposition at a unique stage layer $\Delta^{(s)}$ (Proposition~\ref{prop:same-stage}). The two cases in the definition ($s = m$ vs.\ $s < m$) are mutually exclusive.
\end{proof}

\begin{proposition}[Inclusion compatibility]\label{prop:inclusion-compat}
The inclusion $\iota: W_m \hookrightarrow W_{m+1}$ is an isometry: $\langle \iota(x), \iota(y) \rangle_{W_{m+1}} = \langle x, y \rangle_{W_m}$ for all $x, y \in W_m$.
\end{proposition}

\begin{proof}
$\iota$ embeds $W_m$ as $\bigwedge^1 W_m$, which is orthogonal to all other graded components. The inner product on $\bigwedge^1 W_m$ coincides with that on~$W_m$.
\end{proof}

\begin{proposition}[Metric compatibility]
For $v, w \in \mathcal{H}_m \subset \mathcal{H}_{m+1}$:
\[
  d_{m+1}(v,w) = d_m(v,w),
\]
where $d_m(v,w) = \|\varphi_m(v) - \varphi_m(w)\|_{W_m}$.
\end{proposition}

\begin{proof}
$\varphi_{m+1}(v) = \iota(\varphi_m(v))$ and $\|\iota(x) - \iota(y)\| = \|x - y\|$ by Proposition~\ref{prop:inclusion-compat}.
\end{proof}

Propositions~\ref{prop:dim-tower}--\ref{prop:inclusion-compat} and the metric compatibility above establish that the tower is internally consistent: dimensions grow predictably, embeddings are unique, inclusions are isometric, and distances are preserved across stages. The next three propositions verify that the tower is compatible with the algebraic structure of \S\ref{sec:topology}.

\begin{proposition}[Alternating sign compatibility]
At every level of the tower, permuting the ref list of a $k$-simplex acts on $\Phi_m(\sigma)$ by the sign of the permutation:
\[
  \Phi_m(\sigma \circ \tau) = \operatorname{sgn}(\tau) \cdot \Phi_m(\sigma).
\]
This is the same alternating structure as~$\partial_k$.
\end{proposition}

\begin{proof}
The wedge product is alternating.
\end{proof}

\begin{proposition}[Gram shortcut at every stage]
For $\sigma \in \Delta_k^{(m)}$ with edge vectors $\eta_i = \varphi_m(v_i) - \varphi_m(v_0)$:
\[
  \bigl(\operatorname{vol}_k^{(m)}(\sigma)\bigr)^2 = \det G, \quad G_{ij} = \langle \eta_i, \eta_j \rangle_{W_m}.
\]
This requires $O(k^2 \dim W_m)$ time instead of $O\bigl(\binom{\dim W_m}{k}\bigr)$.
\end{proposition}

\begin{proof}
Standard identity: $\|\eta_1 \wedge \cdots \wedge \eta_k\|^2 = \det(\langle \eta_i, \eta_j \rangle)$.
\end{proof}

\begin{theorem}[Stage enrichment invariance]
If $E$ is an enrichment (modifies only auxiliary keys, preserves $h_{\mathrm{id}}$), then $\Phi_m(\sigma; E(\Sigma)) = \Phi_m(\sigma; \Sigma)$ for all~$m$ and~$\sigma$.
\end{theorem}

\begin{proof}
Enrichment does not change $\mathrm{ref}$ or $h_{\mathrm{emb}}$. The exterior embedding depends only on vertex embeddings (via~$\varphi_m$) and the ref structure. Both are invariant.
\end{proof}

\medskip

So far the tower embeds each stage layer into its own vector space, and the properties above guarantee that this embedding is consistent, isometric, and algebraically compatible. But something new happens at stage $m \geq 1$ that has no analogue at stage~0: the vertex set $\mathcal{H}_{m+1}$ contains entries of \emph{different degrees}. An atom embeds into $\bigwedge^1 W_m$; a $p$-simplex embeds into $\bigwedge^p W_m$. When a simplex in $\Delta^{(m+1)}$ simultaneously references vertices of different degrees, the difference vectors live in different graded components of $W_{m+1}$. This \emph{grade mixing} is the geometric manifestation of the stage decomposition's key feature: meta-knowledge relates objects of fundamentally different structural types. We now make this precise.

\subsubsection*{Grade Mixing}

The vertex set $\mathcal{H}_{m+1}$ contains entries of different degrees. An atom embeds into $\bigwedge^1 W_m$; a $p$-simplex embeds into $\bigwedge^p W_m$. When a simplex in $\Delta^{(m+1)}$ simultaneously references vertices of different degrees, the difference vectors are \emph{mixed-grade} elements of~$W_{m+1}$. This is algebraically well-defined ($W_{m+1}$ is a vector space) and geometrically meaningful.

\begin{definition}[Grade distance spectrum]
For $v, w \in \mathcal{H}_{m+1}$, the difference $\delta = \varphi_{m+1}(w) - \varphi_{m+1}(v)$ decomposes into grade components:
\[
  \delta = \delta_0 + \delta_1 + \cdots + \delta_N, \quad \delta_p \in \textstyle\bigwedge^p W_m.
\]
Since different grades are orthogonal:
\[
  d_{m+1}(v,w)^2 = \|\delta\|^2 = \sum_{p=0}^{N} \|\delta_p\|^2.
\]
The \emph{grade distance spectrum} of the pair $(v,w)$ is
\[
  \mathbf{d}(v,w) = (\|\delta_0\|, \|\delta_1\|, \ldots, \|\delta_N\|) \in \mathbb{R}^{N+1}.
\]
\end{definition}

The grade distance spectrum decomposes a single distance measurement into contributions from each structural level. Two vertices that differ only in their grade~1 components are ``pointing in different directions''; two that differ in grade~2 components have different ``areas''; and so on. The following theorem shows that when the two vertices have different degrees, this decomposition has an irreducible component.

\begin{theorem}[Grade distance decomposition]\label{thm:grade-decomp}
Let $v \in \mathcal{H}_{m+1}$ have degree~$p$ and $w \in \mathcal{H}_{m+1}$ have degree~$q$, with $p < q$. Then:
\[
  d_{m+1}(v,w)^2 = \underbrace{\|\delta_1\|^2 + \cdots + \|\delta_p\|^2}_{\text{shared-grade difference}} + \underbrace{\|\delta_{p+1}\|^2 + \cdots + \|\delta_q\|^2}_{\text{intrinsic complexity of } w}.
\]
The second sum comes entirely from the higher-grade components of~$w$. Even if $v$ and $w$ are perfectly aligned on the shared grades, the intrinsic structural complexity of the higher-degree entity contributes irreducibly to the distance.
\end{theorem}

\begin{proof}
The embedding of~$v$ has grade components only up to $\max(1, p)$; all components at grades $> \max(1,p)$ are zero. Hence $\delta_{p+1}, \ldots, \delta_q$ equal the corresponding components of $\varphi_{m+1}(w)$.
\end{proof}

\begin{example}[Same-grade pair]
Let $e_1 = [a,b]$ and $e_2 = [b,c]$ be two 1-simplices in $\Delta^{(0)}$, both serving as vertices in $\Delta^{(1)}$. Their embeddings in $W_1$ lie in $\bigwedge^1 W_0$:
\[
  \delta = \Phi_0(e_2) - \Phi_0(e_1) = (\varphi_0(c) - \varphi_0(b)) - (\varphi_0(b) - \varphi_0(a)) \in \textstyle\bigwedge^1 W_0.
\]
The spectrum has a single nonzero component. The distance between two relations of the same type is measured exactly as the distance between two atoms: pure direction difference.
\end{example}

\begin{example}[Cross-grade pair]\label{ex:cross-grade}
Let $a$ be an atom and $f = [a,b,c]$ a 2-simplex. In $W_1$:
\[
  \delta = \Phi_0(f) - \iota(\varphi_0(a)) \in \textstyle\bigwedge^1 W_0 \oplus \textstyle\bigwedge^2 W_0.
\]
The grade~1 component measures the positional difference between $a$ and the ``center'' of~$f$. The grade~2 component equals the area of~$f$: it is the intrinsic complexity of the 2-simplex, irreducible and independent of the atom's position.

A ``simple'' face (area $\approx 0$, vertices nearly collinear) has distance from its vertex dominated by grade~1. A ``complex'' face (large area, vertices spread) carries additional grade~2 distance: the structural complexity of the face is itself a form of remoteness.
\end{example}


\medskip

The stage embedding tower and grade mixing give each stage layer a complete geometric structure: distances, volumes, and a graded decomposition of both. The question is now: \emph{when does this geometry faithfully represent the topology?} The answer involves the Kronecker pairing (which connects homology and cohomology) and leads to the central result of this section: the embedding obstruction theorem.

% ============================================================
%   §4.4  EMBEDDING OBSTRUCTION
% ============================================================
\subsection{Embedding Obstruction}\label{sec:obstruction}

Before stating the obstruction theorem, we introduce the pairing that connects homology and cohomology.

\begin{definition}[Kronecker pairing]
\label{def:pairing}
The \emph{Kronecker pairing} $\langle \cdot, \cdot \rangle: H^k(\Delta) \times H_k(\Delta) \to \mathbb{Z}$ evaluates a cohomology class on a homology class:
\[
  \langle [f], [\gamma] \rangle = f(\gamma).
\]
\end{definition}

The pairing connects the two perspectives: homology detects ``where the holes are,'' cohomology measures ``how much a function fails to extend across them.''

\begin{example}
Let $[\gamma] \in H_1$ be the circular dependency $A \to B \to C \to A$, and let $[f] \in H^1$ assign review costs to edges. Then $\langle [f], [\gamma] \rangle = f(e_{AB}) + f(e_{BC}) + f(e_{CA}) = 3 + 2 + 5 = 10$ hours: the total effort debt per traversal.
\end{example}

We now state the central theorem of this section. The intuition is simple: a $k$-cycle requires $k$ independent directions to ``spread out.'' If the embedding space has fewer than $k$ dimensions, every $k$-simplex collapses to zero volume: the geometry cannot see the topology. This gives a hard lower bound on embedding dimension.

\begin{theorem}[Per-stage embedding obstruction]\label{thm:per-stage-obstruction}
For each stage~$m$: if $H_k(\Delta^{(m)}) \neq 0$ and $k > \dim W_m$, then $\operatorname{vol}_k^{(m)}(\sigma) = 0$ for every $\sigma \in \Delta_k^{(m)}$. The $k$-dimensional structure collapses to zero volume.
\end{theorem}

\begin{proof}
At stage~$m$, vertex embeddings lie in~$W_m$. The exterior embedding of a $k$-simplex is an element of $\bigwedge^k W_m$. If $k > \dim W_m$, then $\bigwedge^k W_m = 0$: there are no nonzero $k$-vectors in a space of dimension less than~$k$. Hence $\Phi_m(\sigma) = 0$.
\end{proof}

\begin{corollary}[Per-stage dimension bound]
$\dim W_m \geq k_m^*$, where $k_m^* = \max\{k \mid H_k(\Delta^{(m)}) \neq 0\}$ is the homological dimension of stage~$m$.
\end{corollary}

\begin{theorem}[Stage~0 is the binding constraint]\label{thm:binding}
Since
\[
  \dim W_0 = d \ll 2^d = \dim W_1 \ll 2^{2^d} = \dim W_2 \ll \cdots,
\]
the binding constraint comes from stage~0: $d \geq k_0^*$ is the tightest requirement. Higher stages are exponentially less constrained.
\end{theorem}

\begin{proof}
For a finite astrolabe file, $k_m^* \leq |\mathcal{H}_m| - 1 < \infty$. Meanwhile $\dim W_m \geq 2^d$ for all $m \geq 1$. If $d \geq k_0^*$ and $2^d \geq |\mathcal{H}_1| - 1$ (which holds for any practical~$d$), then $\dim W_m \geq k_m^*$ for all $m \geq 1$.
\end{proof}

\begin{theorem}[Dimensional surplus at higher stages]
For every $m \geq 1$, $\dim W_m \geq 2^d \geq 2^{k_0^*}$. Meta-knowledge has exponentially more room to embed than ground-level knowledge.
\end{theorem}

\begin{corollary}[DAG]
If $\Delta(\Sigma)$ is a DAG (no higher simplices), then $k^* \leq 1$ and a $1$-dimensional grounding suffices. Lean Mathlib, whose dependency graph is acyclic by construction, admits a grounding in $\mathbb{R}^1$ (topological sort).
\end{corollary}

\begin{corollary}[Tree]
If $\Delta(\Sigma)$ is a tree, then $k^* = 0$ and any embedding preserves the topology. Legal codes with a hierarchical chapter/section/article structure fall into this category.
\end{corollary}

\begin{definition}[Topological-geometric inconsistency]
Let $d_{\mathrm{top}}^{(m)}(v,v')$ be the shortest path length in the 1-skeleton of $\Delta^{(m)}$, and $d_m(v,v') = \|\varphi_m(v) - \varphi_m(v')\|$ the embedding distance.
\begin{itemize}
\item $d_{\mathrm{top}}$ small, $d_m$ large: a \emph{cross-domain bridge}. Topologically close, semantically distant. Structurally important.
\item $d_{\mathrm{top}}$ large, $d_m$ small: a \emph{missing relation}. Semantically similar, topologically far apart. A candidate for a new edge.
\end{itemize}
\end{definition}

\begin{example}[Missing relation detection]
In a Lean formalization, \texttt{List.length\_map} and \texttt{Array.size\_map} are semantically close ($d_{\mathrm{emb}} \approx 0.1$) but have no direct reference ($d_{\mathrm{top}} = 5$). The inconsistency flags this as a potential missing lemma: a generalization that unifies both.
\end{example}

\begin{remark}[Grounding quality diagnosis]
Given an embedding~$h_{\mathrm{emb}}$, one can construct the Rips complex $\operatorname{Rips}_\epsilon(\mathcal{H}_0)$ at various scales~$\epsilon$ and compare its homology to~$H_k(\Delta)$. Discrepancies indicate that the embedding distorts the topological structure. This provides a quality metric for embeddings that requires no downstream task.
\end{remark}

\begin{definition}[Weighted homology]
Edge lengths define a weighted boundary operator $\partial_k^{\rho}$ by scaling each face map by the volume of the removed face. The resulting weighted homology $H_k^{\rho}(\Delta)$ detects holes with measurable ``size'': a large cycle in $H_1^{\rho}$ spans a wide semantic region; a small cycle is a local phenomenon.
\end{definition}


% ============================================================
%   §4.5  FORMALIZATION OF AI FRONTIER PROBLEMS
% ============================================================
\subsection{Formalization of AI Frontier Problems}\label{sec:ai-frontier}

The obstruction theorem translates topological complexity into a geometric constraint. We now apply this translation to four problems at the frontier of AI research: scaling saturation (when does increasing embedding dimension stop helping?), hallucination (when does geometric proximity mislead?), the distinction between understanding and memorization (when does a model know positions but not relationships?), and chain-of-thought reasoning (when does step-by-step inference outperform direct shortcuts?). Each problem receives a formal characterization in terms of the stage embedding tower; the characterizations are not metaphors but definitions with testable predictions.

The cup product, introduced next, provides the algebraic tool for the first of these. The obstruction theorem uses the Betti numbers $\beta_k = \operatorname{rank} H_k$; the cup product encodes how obstructions at different levels interact.

\begin{definition}[Cup product]
\label{def:cup-product}
The \emph{cup product} $\smile: H^p(\Delta) \times H^q(\Delta) \to H^{p+q}(\Delta)$ makes the cohomology into a graded ring $H^*(\Delta) = \bigoplus_{k \geq 0} H^k(\Delta)$.
\end{definition}

The cohomology ring $H^*(\Delta)$ is the complete algebraic fingerprint of a knowledge base: it encodes not only the individual obstructions ($H^k$) but also how obstructions at different levels interact. Two knowledge bases with the same Betti numbers but different cup products have qualitatively different dependency structures.

\medskip

\textbf{Scaling saturation (per-stage).} If $\dim W_m < k_m^*$, the embedding dimension at stage~$m$ is insufficient:
\[
  \operatorname{err}^{(m)}(\dim W_m) \geq c \cdot \sum_{k > \dim W_m} \operatorname{rank} H_k(\Delta^{(m)}) > 0.
\]
Stage~0 inflection point: $d = k_0^*$. Stage~1 inflection: $\dim W_1 = k_1^*$, i.e.\ $d = \lceil \log_2 k_1^* \rceil$. Increasing the base dimension~$d$ has \emph{linear} payoff at stage~0 but \emph{exponential} payoff at stage~1: the embedding dimension grows as~$2^d$ there.

\medskip

\textbf{Hallucination (per-stage, grade-stratified).}

\begin{definition}[Stage hallucination set]
\[
  \operatorname{Hal}^{(m)}(\varphi_m, \epsilon, r) = \{(v,w) \in \mathcal{H}_m \times \mathcal{H}_m \mid d_m(v,w) < \epsilon, \; d_{\mathrm{top}}^{(m)}(v,w) > r\}.
\]
\end{definition}

\begin{definition}[Grade-stratified hallucination]
\[
  \operatorname{Hal}^{(m)}_{p,q} = \{(v,w) \in \operatorname{Hal}^{(m)} \mid \deg(v) = p, \; \deg(w) = q\}.
\]
\begin{itemize}
\item $\operatorname{Hal}^{(1)}_{0,0}$: two atoms close in~$W_1$ but topologically far. Same as stage~0 hallucination.
\item $\operatorname{Hal}^{(1)}_{0,1}$: an atom and an edge close in~$W_1$. \emph{Type confusion}: the model conflates an entity with a relation.
\item $\operatorname{Hal}^{(1)}_{1,2}$: an edge and a face close in~$W_1$. \emph{Structural flattening}: the model compresses a ternary interaction into a binary one.
\end{itemize}
\end{definition}

Traditional hallucination analysis captures only $\operatorname{Hal}^{(0)}_{0,0}$. The grade stratification reveals a family of structurally distinct hallucination modes.

\begin{example}[Biomedical type confusion]
In a biomedical knowledge base, a drug entity~$a$ and a drug-interaction edge~$e = [a, b]$ may be close in $W_1$ (the embedding model encodes both as ``related to drug $a$''). A model queried about~$a$ may hallucinate properties that belong to the \emph{interaction} $e$, not to the drug itself. The pair $(a, e) \in \operatorname{Hal}^{(1)}_{0,1}$ is a certificate of this risk.
\end{example}

\medskip

\textbf{Understanding vs.\ memorization (grade version).}

Memorization is a grade-1 phenomenon: a model stores each atom as a point in~$\mathbb{R}^d$. Understanding requires representing relationships (1-simplices, 2-simplices, and higher). A model can perfectly memorize all atoms while completely failing to represent the edges connecting them: the 1-simplices have zero volume.

\begin{definition}[Understanding measure]
\[
  \operatorname{und}(\sigma) = \frac{\operatorname{vol}_k^{(m)}(\sigma)}{\bigl(\max_i \|\varphi_m(v_i)\|\bigr)^k}.
\]
\end{definition}

The normalized $k$-volume. $\operatorname{und}(\sigma) = 0$ means the model knows the positions of all vertices but cannot represent their higher-order independence: it is memorizing, not understanding.

\medskip

\textbf{Chain-of-thought (cross-stage).}

Chain-of-thought reasoning outperforms direct inference precisely on the hallucination set: when two records are semantically close but topologically distant, a direct embedding shortcut fails, but a step-by-step path through the 1-skeleton respects the topological structure.

The stage embedding tower adds a new dimension: CoT can cross stages.

\begin{definition}[Stage-crossing chain of thought]
A \emph{stage-crossing chain of thought} is a path
\[
  v_0 \to \sigma_1 \to \sigma_2 \to \cdots \to \sigma_n \to v_n
\]
where $v_0, v_n \in \mathcal{H}_0$ and the intermediate steps may pass through $\Delta^{(m)}$ for $m \geq 1$.
\end{definition}

``Abstract then reason'' $=$ ascend to a higher stage and walk there. This explains why CoT is especially effective for tasks requiring abstract reasoning: at higher stages, the dimensional surplus (Theorem~\ref{thm:binding}) means the embedding space is exponentially larger, and geometric shortcuts are exponentially less likely to fail.
